\documentclass[UTF8,a4paper,12pt,fontset=fandol]{ctexart}
\usepackage[margin=2.5cm,headheight=0pt,headsep=0pt]{geometry}
\usepackage{amsmath,amssymb,amsthm,mathtools,booktabs,array,tabularx,longtable,graphicx,float,enumitem,xcolor}
\usepackage[font=small,labelfont=bf,labelsep=quad]{caption}
\usepackage[hidelinks]{hyperref}
\usepackage{fancyhdr,fvextra}
\IfFileExists{src/jammer_solver/solver.py}{\def\RepoRoot{.}\def\PaperDir{paper}}{\def\RepoRoot{..}\def\PaperDir{.}}
\pagestyle{fancy}\fancyhf{}\fancyfoot[C]{\thepage}\renewcommand{\headrulewidth}{0pt}
\setlength{\headheight}{0pt}\setlength{\headsep}{0pt}
\linespread{1.16}\setlength{\parindent}{2em}\setlength{\parskip}{2pt}
\setlength{\emergencystretch}{2em}
\setlist{nosep,leftmargin=2em,topsep=4pt}
\ctexset{section={format=\Large\heiti\bfseries,beforeskip=8pt,afterskip=7pt},subsection={format=\normalsize\heiti\bfseries,beforeskip=7pt,afterskip=4pt}}
\newtheorem{proposition}{命题}[section]
\newtheorem{theorem}[proposition]{定理}
\newcommand{\R}{\mathbb R}
\newcommand{\norm}[1]{\left\lVert#1\right\rVert}
\newcommand{\diam}{\operatorname{diam}}
\newcommand{\conv}{\operatorname{conv}}
\newcommand{\dist}{\operatorname{dist}}
\newcommand{\blank}{\underline{\hspace{1.4cm}}}
\newcommand{\datagap}{\textbf{待实测：}\,}
\newcommand{\fig}[3]{\begin{figure}[H]\centering\includegraphics[width=#2\linewidth]{\PaperDir/figures/#1.pdf}\caption{#3}\end{figure}}
\newcommand{\nextpage}{\clearpage}
\newcolumntype{Y}{>{\raggedright\arraybackslash}X}
\graphicspath{{figures/}}
\begin{document}
\begin{center}
{\zihao{2}\heiti\bfseries 基于连续覆盖证书与观测共享调度的\\无线电干扰源自动定位与清除}\par
\vspace{12pt}{\large\heiti\bfseries 摘\quad 要}
\end{center}
针对干扰源个数、位置、接收半径及定向方向未知，且测角误差有界、同地点误差固定的自动清除任务，建立以相容位置集合为信息状态、以实际任务时间为目标的搜索定位模型。将“如何证明不遗漏”与“如何减少移动和检测费用”分别处理，再通过证据账本与预算检查连接，形成可在线执行的完整策略。

针对问题一，将示向度转化为前向楔形的半平面约束，先判定交会区域的空集、无界及有界情形，再通过凸包和旋转卡壳计算直径。构造可由三次测向产生的正三角形反例，证明以区域直径为直径的圆不一定覆盖区域，并给出最远点对中点圆的充要检验条件。

针对问题二，利用“首次已接收”建立源距离与未知接收半径的耦合约束，推导三个圆盘与前向半平面的交作为保证接收候选域。给定移动预算后，将第二测点选择化为可行角区间上的单变量优化，用区间二分、乐观值排序和剪枝缩小两条误差带交集的解析直径上界，再将局部位移还原为全局坐标。

针对问题三，采用原点与半径1125米六点环构成的七站布局。联合维护未知、已发现、已清除及已排除频道，把剩余搜索站和已知源放入同一服务任务排序。在局部定位中复用完整观测历史，在可清除时直接投影到最近的保证清除位置，并通过有限试探额度与下降秩保证备用过程可接续执行。

针对问题四，采用21站布局，并以“合法源位置位于1000米内测点凸包中”为定向发现的充分条件。以连续区域证书核查覆盖；用真实正、负观测估计补测位置的接收概率，选择共享观察点与原地补测。对狭窄可行域构造完整光学覆盖链，通过Voronoi分区逐块证明覆盖后，按预计首次成功成本安排清除次序。

30场离线合成案例中，390个源全部清除。问题三、四各10场常规输入的平均总时间分别为3060.13秒和6261.69秒，移动时间分别占约76.01\%和74.01\%。文中区分离线结果、截图报告的演练成绩与尚未完成的正式测试，正式测试及未测敏感性数据保留空栏。

\vspace{6pt}\noindent\textbf{关键词：}有界误差；交会定位；连续覆盖；观测共享；联合调度；光学清除

\nextpage
\section{问题重述与分析}
\subsection{背景与任务边界}
无线电干扰源的初步监测只能确定目标分布区域，精准定位仍需要近距离观测和现场清除。题目给定半径1800米的圆形目标区域，机器狗从圆心出发，携带测向机、光学探测仪和激光清除设备，在不知道实际源数和位置的条件下自主决策\cite{problem}。任务目标不仅是找到若干个源，还要给出足以结束搜索的证据，并在全部清除的前提下缩短总任务时间。

20个频道中每个频道至多对应一个源，实际源总数为10至16个。全向源在全部方向发射，定向源仅在未知的闭180度角域内发射；有效接收半径分别固定于1000至1500米之间。不同源互不干扰，因此频道提供了可识别的目标身份，但不能提供“该频道是否有源”的先验答案。

机器狗速度为5米/秒，每次测向耗时5秒，改变测向频道另耗时1秒。光学探测耗时3秒，距离源不超过20米时再耗时2秒清除。移动与检测不能同时进行。这些条件使距离最短、观测次数最少与总时间最短成为不同的目标，须在同一个代价模型中权衡。

\subsection{四个问题的输入与输出}
\begin{table}[H]\centering\caption{四问任务分解}\small
\begin{tabularx}{\linewidth}{cYYY}\toprule
问题 & 已知信息 & 主要输出 & 验证重点\\\midrule
一 & 检测点和示向度 & 交会区域直径算法、同直径圆覆盖结论 & 退化状态、最远点对及反例\\
二 & 全向源的一次有效示向 & 第二测点候选域及选择策略 & 保证接收、定位精度与移动量\\
三 & 全部源为全向，具体数目未知 & 搜索、定位、清除及停止策略 & 全清比例、平均定位清除时间\\
四 & 全向与定向混合，类型与方向未知 & 对不可见性稳健的完整策略 & 定向不漏检、成本波动与时间\\\bottomrule
\end{tabularx}\end{table}

前两问研究单源几何信息如何转化为决策，后两问将其嵌入多频道在线任务。问题一回答可行域如何表示；问题二回答下一次观测怎样选择；问题三建立完整搜索与调度；问题四在不可见性增加后修改证据规则和执行方式。四问共享物理参数，但各问的可行域定义不能混用。

\subsection{研究对象与评价指标}
本文以题设协议和本文方法程序为实现对象\cite{interface}。主要指标为全清比例和每个实际清除源的平均任务时间。虚拟任务时间反映移动及设备操作，现实运行时间反映程序计算和通信；两者分别记录。后文的合成模拟、用户提供的演练汇总和正式测试表互相独立，避免用不同来源的数据拼成一次完整成绩。

\nextpage
\subsection{三个相互耦合的困难}
\textbf{第一，误差无法通过同点重复采样消除。} 题设说明同一地点的电磁误差固定。将反复读数当作独立样本并求平均，会人为缩小置信区间。本文保留每次测角对应的整个角区间，通过不同位置的约束相交压缩候选源域。

\textbf{第二，无信号不是无源。} 对全向源，无信号还可能表示超出半径；对定向源，还可能落在发射背面。一个频道只有在全部合法源状态都与实际观测矛盾时，才能从搜索任务中删除。尤其不能因为已经清除10个源，就停止检查其他频道。

\textbf{第三，定位精度与任务费用不等价。} 机器狗可能已经处于能覆盖整个可行域的清除点，却继续为缩小几何直径而移动测量；也可能逐源采用看似最优的定位方案，却让源间往返距离增加。因此，单源动作与多源访问顺序需要同时考虑预计服务出口。

\fig{flow}{0.99}{统一决策流程。实际响应进入证据账本；几何证书和数量上界约束完成判断；费用评分用于安排可执行动作。箭头表示数据或执行结果的传递。}

\subsection{模型之间的连接}
整体流程分为四步。首先，将接收、方向和清除结果转化为集合约束，维护每个频道的状态。其次，用搜索布局或历史证书保证未知源仍有被发现的途径。再次，在剩余搜索与已知源之间选择下一项任务，并在任务内部比较射频测向和光学清除。最后，对动作的预算、实际返回及完成条件做检查，再进入下一轮。

本方法的核心结果是：\textbf{启发式评分可以影响访问次序，但不直接消除一个可能存在的源。} 即使概率估计不准，几何责任区域、实际清除响应和完整备用策略仍决定能否宣告完成。因而可以优化平均费用，同时满足覆盖和实际清除的完成条件。

\nextpage
\section{题设条件、建模约定与符号}
\subsection{采用的条件与补充约定}
\begin{enumerate}
\item 一次案例中，源的位置、频道、有效半径与定向方向固定；不同频道的证据分别维护。该约定对应题目物理环境，不假设源在运行中迁移。
\item 采用平面欧氏距离，机器狗可沿直线移动至目标圆域外。圆域限制源的位置，并不限制机器狗位置，故问题四外环站可位于1800米圆外\cite{interface}。
\item 测角误差按任意有界扰动处理，不假定不同检测位置的误差独立、正态或均匀。理论原始误差为$1^\circ$，程序考虑两位小数舍入后采用$\varepsilon=1.005^\circ$。
\item 证明与执行分开：几何计算采用有理数和向外包络，保证清除判据采用19.8米半径；实际接口的成功阈值仍是20米。
\item 接收概率评分中使用均匀半径和朝向、全向与定向各半的工作假设。这是行动排序的近似模型，题设未给出这些分布，不将其作为排除证据。
\end{enumerate}

\subsection{主要符号}
\begin{table}[H]\centering\caption{跨问题共用的符号}\small
\begin{tabularx}{\linewidth}{p{.24\linewidth}Yc}\toprule
符号 & 含义 & 单位\\\midrule
$\mathcal D,x_c$ & 源分布圆域、频道$c$的源位置 & m\\
$N,R_c,n_c$ & 源总数、接收半径、定向单位法向量 & 个、m、无量纲\\
$p,h$ & 机器狗当前位置与当前测向频道 & m、频道编号\\
$u(\theta),v(\theta)$ & 方位单位向量及其垂直单位向量 & 无量纲\\
$P_c,\widehat P_c$ & 精确相容位置集合及其保守外包 & m坐标集合\\
$U,F,C,A$ & 未知、已发现未清除、已清除、已排除频道集 & 集合\\
$Q_3,Q_4$ & 全向及混合场景的预设搜索站集 & m坐标集合\\
$M,S,K_s,K_f$ & 测向、切频、光学成功及失败次数 & 次\\
$L,T_v,T_r$ & 总行进距离、虚拟时间、现实运行时间 & m、s、s\\
$E,\Phi$ & 剩余试探额度与任务下降秩 & 非负整数\\
$\overline T,\eta$ & 平均定位清除时间与实际全清比例 & s/源、无量纲\\\bottomrule
\end{tabularx}\end{table}

所有角度在文字中用度表示，在三角函数内部换算为弧度。局部二维叉积记作$[a,b]=a_xb_y-a_yb_x$。不同问题的辅助变量在引入处定义，避免把可行域直径、接收半径和清除半径都记成同一个字母。

\nextpage
\section{统一观测模型与时间目标}
\subsection{可接收、方向与近距观测}
令$\mathcal D=\{x\in\R^2:\norm{x}\le1800\}$。全向源在位置$q$可被接收，当且仅当$\norm{q-x_c}\le R_c$。定向源还须满足
\begin{equation}\label{eq:detect}
 n_c^\mathsf T(q-x_c)\ge0,\qquad \norm{n_c}=1.
\end{equation}
这里$n_c$是从源指向有效发射半平面的法向量，不能与检测点指向源的测向向量混淆。

若可接收且距离大于5米，响应为方向值$\theta$，真实源落在以$q$为顶点、方向区间$[\theta-\varepsilon,\theta+\varepsilon]$为边界的前向楔形内。若可接收且距离不超过5米，返回近距响应。若不可接收，返回无信号。近距响应直接给出5米圆盘，不应人为补出一个方向值。

\subsection{动作费用与频道状态}
设前后动作位置为$p,q$，当前测向频道为$h$。一次测向与一次光学动作的费用分别为
\begin{align}
 t_m(p,q,c,h)&=\frac{\norm{q-p}}5+5+\mathbf1_{c\ne h},\\
 t_o(p,q,c)&=\frac{\norm{q-p}}5+3+2\mathbf1_{\text{清除成功}}.
\end{align}
光学动作指定目标频道，但不改变测向机的当前频道，也不产生切频费用\cite{interface}。进入和主动退出不推进虚拟时钟。于是整场费用严格分解为
\begin{equation}\label{eq:cost}
 T_v=\frac L5+5M+S+5K_s+3K_f.
\end{equation}
该公式用于逐条响应回放，是结果表和费用图的共同计算依据。

\subsection{优化与完成约束}
令策略$\pi$根据截至当前的真实观测历史$H_t$选择下一动作。理想目标是对所有合法源状态保证$K_s=N$，并使$T_v$尽量小。本文实际采用有界计算预算的在线近似决策，而不求整个连续信念空间的精确最优策略。
\begin{equation}
 \eta=\frac{K_s}{N},\qquad \overline T=\frac{T_v}{K_s}\quad(K_s>0).
\end{equation}
由于接口不直接返回$N$，$\eta$只能在演练结束获得真值后核验。在线停止条件由频道证据定义，不能以预测的$\eta=1$替代。

现实运行时间还包括求解和HTTP通信。算法始终服从进入时返回的剩余窗口；本文的几何完备性结论以有效响应、正确数值包络和足够执行时间为条件，不据虚拟时间较小推断现实窗口必然满足。

\nextpage
\section{问题一：交会定位区域与直径}
\subsection{前向楔形的半平面表达}
令第$i$个检测点为$p_i$，示向度为$\theta_i$，$u(\theta)=(\cos\theta,\sin\theta)$。对应的前向楔形为
\begin{equation}
 W_i=\{x:[u(\theta_i-\varepsilon),x-p_i]\ge0,\ [u(\theta_i+\varepsilon),x-p_i]\le0\}.
\end{equation}
因为$2\varepsilon<180^\circ$，上述两式确定的是朝向测量方向的凸楔形。若只保留“靠近示向直线”的条件，会把检测点背后的对称区域错误地纳入定位区域。

将每个楔形写成两条线性不等式，得到
\begin{equation}\label{eq:halfplanes}
 P=\bigcap_{i=1}^nW_i=\{x:a_j^\mathsf Tx\le b_j,\ j=1,\ldots,2n\}.
\end{equation}
问题一计算的对象是这个纯交会区域。用于后续行动时可以再与源圆域、已接收半径上限等先验相交，但不能在回答问题一时悄悄加入人工边界盒，使原本无界的交集变成一个有限多边形。

\subsection{空、无界与有界状态的分类}
枚举所有非平行边界线对。若$a_j,a_k$线性独立，其交点通过二元线性方程求出；只有满足全部$2n$条不等式的点才是可行候选。存在至少一个尖楔形时，非空交集不含整条直线，因此存在极点。精确算术下，找不到可行候选可以判空；浮点实现对近平行数值情况保留未决标记。

对非空交集，再考察衰退锥
\begin{equation}
 \mathcal R=\{d:a_j^\mathsf Td\le0,\ j=1,\ldots,2n\}.
\end{equation}
若存在非零$d\in\mathcal R$，则对可行点$x_0$有$x_0+td\in P$，$t\ge0$，故直径无穷。二维尖锥的极射线与某约束边界平行，可枚举约束法向量的两个垂直方向并检验。没有观测时直接返回全平面无界状态。

\subsection{实现中的数值判据}
纯交会模块采用浮点边界枚举，对可行性使用相对容差；近平行交点、无穷值及空集按状态区分返回。后两问的证书模块另外使用有理数裁剪。本文不把前者的浮点直径输出解释为严格区间证明；这一区别也决定了问题一结果的保留小数位和后续行动的安全余量。

\nextpage
\subsection{直径计算与复杂度}
\begin{proposition}\label{prop:diam}
若$P$是非空有界凸多边形，顶点为$z_1,\ldots,z_m$，则
\begin{equation}
 \diam(P)=\max_{i,j}\norm{z_i-z_j}.
\end{equation}
\end{proposition}
\begin{proof}
任意$x\in P$可写成顶点的凸组合$x=\sum_i\lambda_i z_i$。对固定$y$，三角不等式给出$\norm{x-y}\le\sum_i\lambda_i\norm{z_i-y}\le\max_i\norm{z_i-y}$。再对$y$作凸组合展开，即可把最大值同时推到两个顶点上。
\end{proof}

对候选点求逆时针凸包后，采用旋转卡壳枚举对踵点对\cite{toussaint}。边界交点枚举有$O(n^2)$对，每点检验$O(n)$个条件，因此整体为$O(n^3)$；顶点凸包为$O(m\log m)$，卡壳扫描为$O(m)$。在本题小规模观测下，直接枚举易于检查，也能保留清楚的退化状态。

\begin{table}[H]\centering\caption{交会区域直径算法}\small
\begin{tabularx}{\linewidth}{cY}\toprule
步骤 & 运算与输出\\\midrule
1 & 为每次示向生成前向楔形的两条半平面约束。\\
2 & 枚举非平行边界交点，检验全部约束，保留可行点。\\
3 & 若无可行点，返回空或数值未决；若存在非零衰退方向，返回无界。\\
4 & 对剩余点求凸包；单点直径为0，线段直径为端点距离。\\
5 & 对多边形执行旋转卡壳，返回最大距离和对应最远点对。\\
6 & 用最远点对中点及所有顶点检验同直径圆的覆盖性。\\\bottomrule
\end{tabularx}\end{table}

\subsection{同直径圆的充要检验}
设最远点对为$a,b$，$D=\norm{a-b}$。任何半径$D/2$且覆盖$a,b$的圆，其圆心必须是$c=(a+b)/2$。因此存在同直径覆盖圆的充要条件是
\begin{equation}\label{eq:diametercircle}
 \max_j\norm{z_j-c}\le D/2.
\end{equation}
只检验这两个最远点不够；所有其他顶点也必须满足条件。若条件不成立，可另求最小包围圆或使用更保守的半径，但不能把$D/2$直接作为保证覆盖半径。该结论为后续“整个可行域均在20米内”的清除判据提供了必要修正。

\nextpage
\subsection{可实现的反例与算例}
考虑边长为$d$的正三角形。其区域直径为$d$，最小包围圆半径却为$d/\sqrt3$，大于$d/2$。故问题一的第二问答案为\textbf{不一定能覆盖}，且失败并非浮点误差造成。

这一反例可以由符合测向结构的楔形实际产生。沿单位正三角形各条逆时针边向后延长100米设置检测点，将示向度设为该边方向再加$1^\circ$，误差半宽取$1^\circ$。每个楔形的一条边界与三角形边重合，另一条边界不切入三角形；三个前向楔形之交恰为该三角形。这样，反例既满足凸几何性质，也对应题目允许的观测形式。

\fig{q1_geometry}{0.99}{问题一的数值交会与解析反例。左图取检测点$(0,0)$、$(100,0)$及示向度$45^\circ$、$135^\circ$，误差半宽$1.005^\circ$；虚线连接最远点。右图以边长$d$归一化，虚线半径为$d/2$，点划线为最小包围圆。}

\begin{table}[H]\centering\caption{问题一的确定性检验结果}\small
\begin{tabularx}{\linewidth}{YYY}\toprule
输入情形 & 输出或理论结果 & 核验内容\\\midrule
单个正常示向 & 无界楔形 & 不引入人工有限边界\\
前向条件相互矛盾 & 空或数值未决 & 不输出虚构有限直径\\
图中两次示向 & $D=3.5095516$米 & 顶点与最远点对一致\\
单位正三角形 & $D=1$，$r_{\min}=1/\sqrt3$ & 同直径圆覆盖判据不成立\\\bottomrule
\end{tabularx}\end{table}

在实际多次测向中，新增有效约束使精确交集满足$P_{k+1}\subseteq P_k$，故直径不会增大。这一集合嵌套性比假设平均误差按样本量缩小更适合本题。若计算结果反而扩张，首先应检查角度环绕、前向方向、单位换算与数值容差，而不能把扩张当作真实源运动。

\nextpage
\section{问题二：第二检测点的候选域}
\subsection{首次接收给出的半径耦合}
设首次正常测向位置为$p_1$、示向度为$\theta_1$。在局部正交基$u=u(\theta_1)$、$v=v(\theta_1)$下，真实源可写为
\begin{equation}
 x=p_1+r(\cos\delta\,u+\sin\delta\,v),\quad5<r\le1500,\quad|\delta|\le\varepsilon.
\end{equation}
因为首次已经接收到信号，未知半径并非独立自由参数，而满足$R\ge\max\{1000,r\}$。设第二点为$q=p_1+au+bv$，选择前向半平面$a\ge0$。把$r$放宽到$[0,1500]$后，保证接收要求
\begin{equation}\label{eq:q2robust}
 \norm{(a,b)-r(\cos\delta,\sin\delta)}\le\max\{1000,r\},\quad\forall r,\delta.
\end{equation}

\subsection{消去连续距离变量}
当$0\le r\le1000$时，距离平方关于$r$为凸函数，最大值出现在$r=0$或$r=1000$。当$r\ge1000$时，平方后抵消$r^2$，式\eqref{eq:q2robust}等价于
\begin{equation}
 a^2+b^2\le2r(a\cos\delta+b\sin\delta).
\end{equation}
若$r=1000$时成立，右侧系数必非负，增大$r$不会使约束更紧。对$a\ge0$及小角区间，最不利角度是两个区间端点。于是只需检查三个圆盘：
\begin{equation}\label{eq:q2region}
\begin{split}
\mathcal C={}&B((0,0),1000)\cap\{a\ge0\}\\
&\cap B((1000\cos\varepsilon,1000\sin\varepsilon),1000)\\
&\cap B((1000\cos\varepsilon,-1000\sin\varepsilon),1000).
\end{split}
\end{equation}
该集合是放宽先验和前向条件下的完整保证接收域；对真实$r>5$及源圆域约束，它是保守的充分候选域。若首次响应为近距，不需要使用此构造，可直接在原位置清除。

\noindent\textbf{首次接收信息的作用。}
若忽略首次接收事实，要求全部1500米端点都落在第二点的1000米接收圆内，会排除许多合法候选。例如$(a,b)=(500,840)$满足式\eqref{eq:q2region}，却不满足上述独立处理。其改进来自保留观测造成的变量耦合，而不是放松真实接收约束。

\nextpage
\subsection{定位精度的解析上界}
候选域保证能够接收，但不自动保证交会角度良好。令$\rho=\sqrt{a^2+b^2}$、$\alpha=\operatorname{atan2}(b,a)$，先考虑$b>0$的一侧。真实源到第二点的距离上界为
\begin{equation}
 d_2(\rho,\alpha)=\max_{r\in\{5,1500\}}\sqrt{r^2+\rho^2-2r\rho\cos(\alpha+\varepsilon)}.
\end{equation}
第二点相对于首次方向的最小横向高度为$h=b\cos\varepsilon-a\sin\varepsilon$。若$h>0$，两条名义测向方向的最小交角可保守估计为
\begin{equation}
 \gamma=\arcsin\!\left(\min\{1,h/d_2\}\right)-2\varepsilon.
\end{equation}
当$\gamma\le0$时，无法由此得到有限的条带交直径上界。令两条包络条带的半宽为$w_1=1500\sin\varepsilon$、$w_2=d_2\sin\varepsilon$，则最宽对角线给出
\begin{equation}\label{eq:q2bound}
 D_{\mathrm{ub}}=\frac{2\sqrt{w_1^2+w_2^2+2w_1w_2\cos\gamma}}{\sin\gamma}.
\end{equation}
这是对物理相容区域的外包估计。缩小该上界有助于获得稳健交会，但它与真实区域直径的最小化不是同一问题。

\subsection{给定移动预算的第二测点选择算法}
输入为首次测点$p_1=(x_1,y_1)$、示向度$\theta_1$和移动预算$\rho$，其中$5<\rho<1000$米。取第二点到首次测点的距离为$\rho$，只优化移动方向。输出为第二测点$q^*$、对应的解析直径上界及实际计算间隙。

\noindent\textbf{步骤1：建立局部坐标与可行角区间。}
以首次示向为前向轴$u=(\cos\theta_1,\sin\theta_1)$，取垂直轴$v=(-\sin\theta_1,\cos\theta_1)$。先在$b>0$侧搜索，令$a=\rho\cos\alpha$、$b=\rho\sin\alpha$。将其代入式\eqref{eq:q2region}，结合横向高度$h>0$，得到
\begin{equation}\label{eq:q2angles}
 I_0=[\varepsilon+\delta_\alpha,\ \arccos(\rho/2000)-\varepsilon-\delta_\alpha],
 \qquad\delta_\alpha=10^{-8}\ \mathrm{rad}.
\end{equation}
全部角度在计算前转换为弧度；$\varepsilon=1.005^\circ$。端点内缩用于保留数值余量。$\rho$相对于首次测点计量，给定当前机器狗位置不会改变此约束。

\noindent\textbf{步骤2：评价候选角度并初始化。}
对任一$\alpha\in I_0$，计算前述$d_2$、$h=\rho\sin(\alpha-\varepsilon)$和$\gamma$，再由式\eqref{eq:q2bound}得到目标值$F(\alpha)=D_{\mathrm{ub}}(\rho,\alpha)$。若$\gamma\le0$，令$F(\alpha)=+\infty$，表示该角度不能由此解析式给出有限上界。先评价$I_0$的两端及中点，保存最小值$U$和对应角度$\alpha^*$，再将整个区间放入优先队列。

\nextpage
\noindent\textbf{步骤3：估计每个区间可能达到的最好值。}
记$\Phi(d,h)$为按上节交角和条带公式得到的直径上界。对区间$I=[\alpha_l,\alpha_r]$，$d_2$随角度增大而增大，$h$也随角度增大而增大。取区间中较小的距离和较大的横向高度，构造乐观值
\begin{equation}\label{eq:q2interval}
\begin{split}
 d_{\min}(I)&=d_2(\rho,\alpha_l),\qquad
 h_{\max}(I)=\rho\sin(\alpha_r-\varepsilon),\\
 L(I)&=\max\{0,\ \Phi(d_{\min}(I),h_{\max}(I))-10^{-6}\ \mathrm m\}.
\end{split}
\end{equation}
这两个极值不要求在同一角度同时取得，因而用于估计该区间还有多少改进空间。减去$10^{-6}$米作为浮点计算余量。优先队列按$L(I)$从小到大排列，使尚有较大改进可能的角区间先被处理。

\noindent\textbf{步骤4：二分、更新和剪枝。}
每次取出队首区间$I=[\alpha_l,\alpha_r]$。若$U-L(I)\le\tau$，停止搜索；否则在中点将其分成左右两个子区间。在每个子区间的中点计算$F$，若小于$U$，便更新$U$和$\alpha^*$。随后计算子区间的$L$：当$L<U$时放回队列；当$L\ge U$时，该区间不能改善当前结果，将其剪除。被剪除区间的最小乐观值仍参与最终间隙计算。

\noindent\textbf{步骤5：处理停止条件并报告实际间隙。}
默认$\tau=0.02$米，该量是直径目标值的容差；每评价一个子区间中点，计数增加1，评价上限为1024次，初始化的三次评价不计入此数。队列为空、队首满足容差或评价数达到上限时结束。综合当前可行值、队列内区间、被剪除区间及触发停止区间的乐观值，得到$L_{\mathrm{all}}$，输出
\begin{equation}
 G=U-L_{\mathrm{all}}.
\end{equation}
达到评价上限时，$G$可能仍大于0.02米，应报告实际数值；这不是整个物理定位问题的最优性间隙。若结束后仍没有有限的$U$，程序提示该移动预算下未得到有限解析上界，需要增大移动预算。

\noindent\textbf{步骤6：还原全局坐标并选择对称侧。}
由最终角度计算$a^*=\rho\cos\alpha^*$、$b^*=\rho\sin\alpha^*$。候选域关于前向轴对称，因此得到两个可用位置
\begin{equation}\label{eq:q2positions}
 q_+=p_1+a^*u+b^*v,\qquad q_-=p_1+a^*u-b^*v.
\end{equation}
若输入当前机器狗位置$p_c$，选取$\norm{q-p_c}$较小的点；距离相同时选$q_+$。未给出$p_c$时也选$q_+$。最终返回$q^*$、$U$、$G$及保证接收的距离余量，后续测向仍按实际接口返回更新源的可行区域。

\nextpage
\section{问题三：按运行顺序组织全向源搜索与清除}
问题三的所有干扰源均为全向源。只要机器狗进入某源的有效接收半径，就能收到该频道的信号。因此，算法先用少量搜索点保证每个源都有被发现的机会，再根据实际测向逐步缩小源的位置范围，在能够保证清除时立即清除。整场运行采用\textbf{安排任务、取得观测、更新范围、继续定位或清除、重新安排任务}的循环，把搜索、定位与清除衔接起来。

\subsection{运行开始：建立频道记录与七个搜索点}
机器狗从原点出发，对20个频道分别保存状态与观测记录。初始状态均为“未知”$U$；收到信号后成为“已发现但未清除”$F$；清除成功后成为“已清除”$C$；确定不存在源的频道记为“已排除”$A$。同一个位置可以检查多个频道，但每条返回只更新对应频道，不能相互替代。

用于发现未知源的初始搜索点为原点和半径1125米圆周上等角分布的六点：
\begin{equation}\label{eq:q3stations}
Q_3=\{(0,0)\}\cup\{1125(\cos(k\pi/3),\sin(k\pi/3)):k=0,\ldots,5\}.
\end{equation}
七点布局的依据是题设给出的最小接收半径。每个源的有效半径至少为1000米，所以只需让七个1000米圆盘覆盖整个1800米源域。距原点不超过1000米的源由中心站发现；外侧的源则分配给方位最接近的环站，两者相对原点的方位差不超过$30^\circ$。在这一范围内，源到对应环站的最大距离约为999.11米，仍小于1000米。

\fig{coverage_layouts}{0.99}{初始搜索点布局。左图中，问题三的七个1000米接收圆覆盖1800米源域，使每个全向源至少能在一站被接收。右图为问题四所用的21站布局。搜索点确定可发现的范围，实际访问次序随观测和清除进度调整。}

这说明逐站检查足以发现所有源，并不要求每次都走完七站。若某些频道已经解决，相关测量便可省去；若历史检测位置能够替代后续站点，也可以缩短搜索路线。是否真的可以省去一站，需要检查调整后的点集仍能覆盖全部可能的源位置。

\subsection{每轮开始：把搜索与清除放在一起安排}
程序首先整理当前未完成的任务：一类是去搜索站检查未知频道，另一类是处理已发现但未清除的源。两类任务一起安排，是因为搜索和清除都消耗移动时间。例如，刚发现的源若就在附近，先清除它再去下一站，可能比走完所有搜索站后折返更省时。

每项任务$j$有一个或多个执行方式$m$，用入口$e_{jm}$、预计服务时间$s_{jm}$和预计出口$x_{jm}$表示。搜索任务的入口、出口都是站点，站内时间按需要检查的频道数乘以6秒估计。源任务若已可直接清除，入口就是清除位置；否则提供不同测向入口，并以当前源域中心估计服务结束位置。该中心仅用于估计行程，不当作已经找到的源坐标。

对一组任务顺序$\sigma$和执行方式$m_i$，估计总时间为
\begin{equation}\label{eq:service}
\widehat T=\frac{\norm{p-e_{\sigma_1m_1}}}5+s_{\sigma_1m_1}
+\sum_{i=2}^{|\mathcal J|}\left[\frac{\norm{x_{\sigma_{i-1}m_{i-1}}-e_{\sigma_im_i}}}5+s_{\sigma_im_i}\right],
\end{equation}
其中$p$是当前位置，$\mathcal J$为待处理任务集合。式中把“走到任务入口的时间”和“任务本身的时间”加在一起，因此不会只按两个源域中心之间的距离排序。

程序先比较沿用上轮顺序、就近贪心、先处理源、先访问搜索站，以及通过指派方法连接任务得到的顺序。指派步骤先为任务寻找总连接费用较小的前后关系，再把形成的小环接成一条完整顺序，其基础为匈牙利方法\cite{kuhn}。对固定顺序，用动态规划挑选各任务的执行方式：到达当前任务的累计费用，等于前一任务的累计费用、两者之间的移动费用和当前服务费用之和。

选出较好方案后，再做至多两轮任务移位，尝试把某项任务插入别的位置。最终先执行排在首位的任务。搜索任务结束后根据新发现重新排序；选中一个源后，通常持续处理到清除，再安排其余任务。这样既考虑后续行程，也保留随观测调整计划的余地。

\subsection{到达搜索站：检查未知频道并建立源的位置范围}
机器狗到达选定搜索点后，逐一检查仍未知且尚未在此点测过的频道；优先测当前已调谐的频道，减少切频。随后根据实际返回处理。

\textbf{返回方向值时，建立位置区域。} 测向有误差，源并不一定在示向直线上。算法以检测点为顶点，保留示向两侧完整误差范围内的前向区域，并结合最大接收距离与源域边界，形成包含真实源的保守多边形$\widehat P_c$。计算采用$1.005^\circ$包住题设$1^\circ$误差。后续再收到方向值，就将新观测对应的区域与旧区域求交，而不是直接取两条示向线的交点。

\textbf{返回近距时，准备直接清除。} 近距表示源在当前点5米内，程序用边长10米的方盒保守表示它的位置。整个方盒都在20米清除范围内，随后即可选择直接清除，无须继续测向求精确坐标。

\textbf{没有信号时，继续排查该频道。} 全向源在1000米内必能被接收，因此本次阴性排除了该点1000米圆内存在该频道源的可能。但圆外仍可能有源，频道暂时保持未知。只有该频道积累的实际阴性点足以覆盖整个源域，才能改为已排除；尚未到达的搜索点不算已经取得的观测。

首次发现源时，程序还为其位置区域准备一套光学排查方案：沿首次测向轴每隔25米安排一个纵向位置，横向取$-25,0,25$米三层，最多形成183个小块及对应清除点。各块仅保留与当前区域的交，且整块必须位于对应点的20米圆内。这样，后续即使测向收益有限，也有覆盖剩余位置范围的清除途径。下文将这些尚需处理的小块称为“责任块”。

\subsection{同站发现多个源：用一次移动共享第二次观察}
若同一搜索站至少新发现三个源，且每个源都只有一次方向观测，程序会考虑先到一个共同观察点，对这些频道分别再测一次。目的在于让同一段移动同时服务多个源，避免以后为每个源单独绕行。

候选点位于当前搜索站周围350米的圆上，每隔$30^\circ$取一点，共12个方向。一个候选方向是否有用，主要看它能否改变原来的观察角度。设某源的当前区域沿原测向轴的宽度为$w$，候选点到区域代表位置的距离为$d$，新旧观察方向夹角为$\phi$，补测后的宽度近似为
\begin{equation}\label{eq:q3width}
\widehat w'=\min\left\{w,\frac{2\varepsilon d}{\max(0.01,|\sin\phi|)}\right\},
\qquad \varepsilon=1.005\pi/180.
\end{equation}
两次观察几乎平行时，$|\sin\phi|$很小，交会区域沿长轴仍可能很长；从侧面观察形成较大夹角时，同样的测角误差能给出更窄的距离范围。

程序把宽度减少量的一半近似折算为可节省的移动距离，以5米每秒换算后，每个源的收益估计为$(w-\widehat w')/10$秒。再用总收益与绕行费用比较，对候选点$q$计算
\begin{equation}
S(q)=\frac{350+\norm{q-t}-\norm{p-t}}5
-\sum_c\frac{w_c-\widehat w'_c}{10},
\end{equation}
其中$p$为当前搜索站，$t$取最近的已知源域中心，作为后续去向的参考。选择$S(q)$最小的方向，并在行动额度和时间允许时逐频道补测。

当前问题三的实现用该分数比较12个方向，没有另设“总收益必须为正才移动”的门槛，也没有在此处建立接收概率模型。因此，这一步是对共享位置的启发式选择，实际能节省多少时间仍取决于观测结果。补测无信号时，照常记录并利用全向阴性约束更新范围。

\subsection{选中一个源：先检查能否从某个位置直接清除}
任务排序选中源$c$后，首先判断当前区域是否已经小到可以一次清除。需要找到位置$q$，使源无论在$\widehat P_c$中的哪里，都距$q$不超过19.8米。若当前位置为$p$，先考虑
\begin{equation}\label{eq:clearprojection}
q^*=\arg\min_q\norm{q-p}\quad\text{s.t.}\quad
\norm{q-z}\le19.8,\quad\forall z\in\operatorname{vert}(\widehat P_c).
\end{equation}
多边形是凸的，检查所有顶点就能确认整个区域都在该清除圆内。19.8米比题设20米略小，用于留出数值余量。

求解时，可以把每个顶点看作一个半径19.8米圆的圆心。所有圆的公共部分，就是能够保证清除的位置集合。若当前位置已在其中，原地清除即可；否则，程序比较圆弧上的最近位置和圆与圆的交点，寻找较近的可行位置。若已有下一任务入口$t$，还在可行范围内尝试减小$\norm{p-q}+\norm{q-t}$，兼顾清除前后的两段移动。

这一步使定位服务于清除：只要整个位置区域已能被一次清除覆盖，就不再为获得更精确的坐标继续花费测向时间。若暂时不能直接清除，则按任务排序给出的入口先取得新观测，再进入源内的反复定位与清除决策。

\subsection{还不能直接清除：选择一个新的测向点}
继续定位时，以最近一次确实收到方向值的位置$o$为参考，令$u$为示向单位向量，$v$为垂直方向，$L_c$为从$o$到源的当前距离上界。只有一次方向观测时，在沿轴前进的不同距离上向两侧偏移，形成候选点
\begin{equation}
q_\pm=o+\bigl[5+(L_c-5)\lambda\bigr]u\pm(L_c/50)v,
\qquad\lambda\in\{0.38,0.50,0.62\}.
\end{equation}
沿轴前进使机器狗接近源，侧向偏移增加交会角度。对每个前进比例，程序采用距当前位置较近的一侧作为本轮候选。\textbf{问题三一轮只执行选定的一个测向点}，得到结果后先更新区域，再选择下一步。

获得至少两次方向观测后，已有信息通常已排除最初范围中的大部分位置。此时把当前多边形投影到最新示向轴上，得到剩余区间$[\ell,h]$，取其中心$(\ell+h)/2$作为前进位置，侧向距离改为$\max\{5,(h-\ell)/20\}$米。步长随剩余范围缩小，避免反复从首次观测的完整距离范围重新搜索。

一个源刚被选中时，入口来自整体任务排序。后续源内选择则比较“测完以后还剩多少清除工作”：将可能返回的方向分成六个角区间，估算每个区间下保留的责任块及其排查时间，再加上到测点的移动、测向和切频费用。同时计入近距分支；若尚不能保证该点对整个区域都在1000米内，还计入无信号后继续排查的费用。用较差分支的费用比较候选，避免只追求某一种有利返回。

程序也同时准备直接按责任块排查的方案。当前瞻没有明显优势、但剩余块仍多于四块时，允许再尝试一次测向，以避免过早走很长的光学路线；源内方向更新轮数达到限制或可选额度不足时，转向光学排查。所有前瞻都只是对可能结果的估算，只有实际返回能够修改源的位置区域。

\subsection{取得测向结果：有方向就求交，无信号就利用距离关系}
若本轮返回方向值，就将新误差区域与旧区域求交，并把本次检测点与方向作为新的参考；若返回近距，则缩小到近距范围。更新后重新检查是否已能直接清除。

全向源的阴性观测同样有用。设$p$曾经收到该频道的信号，$q$本次没有收到，真实源位置为$x$，固定但未知的接收半径为$R$，则
\begin{equation}
\norm{x-p}\le R<\norm{x-q}.
\end{equation}
也就是说，源一定比起$q$更靠近$p$。把两边距离平方并消去$\norm{x}^2$，得到
\begin{equation}\label{eq:omnineg}
2(q-p)^\mathsf T x<\norm q^2-\norm p^2.
\end{equation}
这是一条半平面约束：源只能在两测点垂直平分线靠近阳性点的一侧。例如，$p=(0,0)$、$q=(1000,0)$时，仅这一对观测就要求源的横坐标小于500米。无需知道$R$的具体值，也能排除另一侧的位置。

程序将新阴性分别与已保存的阳性点组合，把所得半平面与当前区域相交。为避免舍入误删边界，计算保留非严格的不等式外包；同时删除完全位于$q$的1000米圆内的责任块，因为全向源若在这些块中，本次必能被接收。整个过程利用的是实际接收与未接收的距离关系，不把重复读数当成新的独立样本。

\subsection{选择光学排查时：先执行一个清除点，再根据结果继续}
若一次清除尚不能覆盖整个区域，而继续测向已不划算或达到轮数限制，算法就使用现有责任块对应的清除点。每个块已保证能从其对应点清除，因此这些点合起来提供完整排查途径。

程序比较两种访问次序：按最初块编号形成的蛇形顺序，以及每次去最近剩余清除点的顺序。评价时把移动时间、前面各次失败探测的3秒和最终成功探测清除的5秒累计，取首次成功可能出现在不同位置时的最大费用。这样比较的是完整排查所需的保守时间，而不是只比较第一段路程。

正常运行先执行较好顺序的第一个清除点。若返回成功，立刻把该频道改为已清除并取消后续动作；若返回未命中，只删除被本次20米圆完整覆盖的责任块，其他块继续保留。随后用剩余块形成保守位置区域，重新判断直接清除、继续测向或下一次光学排查。一次失败只说明本次范围内没有目标，不能把整个频道删除。

这种逐次更新允许先前失败帮助缩小后续工作，同时始终保留真实源可能所在的块。执行固定备用过程时，则按未完成块的顺序继续排查，直到实际清除成功。

\subsection{完成当前任务：原地补测其他源，再安排后续行程}
一次搜索及其共享补测结束后，或一个源清除完成后，机器狗先检查能否利用当前位置补测其他已发现源。既然已经走到这里，原地切频测量无需额外移动；如果视角合适，就可能省去以后的一段绕行。

候选源需要满足：建立源记录后保存的正、负测点合计不足三个；当前点未对该源测过；当前位置距其区域中心不超过1000米。再用式\eqref{eq:q3width}估算新宽度，要求视差$|\sin\phi|\ge0.25$，且$(w-\widehat w')/10-6>0$。按这一净收益从大到小补测。这里的中心距离是筛选依据，未声称整个区域一定可接收；实际阴性仍按式\eqref{eq:omnineg}更新。

随后回到任务排序。已清除源从任务中删除，其余源用新区域重新估计入口和费用。搜索路线也随之调整：先去掉所有剩余未知频道都已检查过的站点，再尝试删去冗余站或用当前清除位置替代某一站。只有共同的历史检测点与调整后的未来站点仍足以覆盖源域，才采用新路线。这里先缩短的是未来行程；频道是否无源仍由真实检测记录决定。

\nextpage
\subsection{行动前检查预算：必要时接续完成备用过程}
每项可选行动之前，算法估计剩余任务沿备用路径完成的费用，再检查额外移动和测量是否仍能放进虚拟时间及现实时间窗口。若没有足够余量，或额外试探额度接近耗尽，就转入备用过程，不把剩余时间继续用于尝试优化。

备用过程先清除已发现源，按剩余责任块逐点排查；再沿间距600米的$7\times7$网格检查未知频道。源域内任一点位于某个网格方格中，到其四角的距离均不超过$600\sqrt2<1000$米，因此网格能继续承担完整搜索。遇到新源就先完成该源的光学排查，再回到网格搜索。已经完成的检测和清除会保留，备用过程接续当前进度。

为使额外尝试不会无限重复，程序维护贯穿整场的有限额度$E$。用$G_c$表示未知频道未检查的备用网格点，$\mathcal B_c$表示已知源剩余责任块，$D_t=F_t\cup C_t$表示已经发现的不同源，可把剩余工作写成
\begin{equation}\label{eq:rank}
\Phi=\sum_{c\in U}|G_c|+\sum_{c\in F}|\mathcal B_c|+183(16-|D_t|)+E.
\end{equation}
前两项计算已经明确的待查任务，第三项预留最多还可能发现的源所带来的责任块。新发现一个源时，这项预留随之转为实际清除任务；检测、排查或清除使任务减少。若一次行动没有使这些任务量严格下降，就消耗一个不可重置的额外额度。因此，算法不会因反复重规划而获得无限试探机会。

在题设物理规则成立、通信有效且执行窗口足够的条件下，保守位置区域、完整责任块和有限的备用过程共同提供完成途径。实际是否完成仍由后面的停止条件判断。

\subsection{整场结束：所有已知源清除且其余频道均被排除}
每轮循环检查
\begin{equation}\label{eq:stop}
U_t=\varnothing,\qquad F_t=\varnothing,\qquad 10\le|C_t|\le16.
\end{equation}
第一项表示不再有尚未解决的频道，第二项表示没有发现后尚未清除的源，第三项要求实际清除数符合题设范围。清除状态只由真实成功响应产生。

未知频道可以在其实际阴性检测点已经覆盖整个源域时被排除；若已发现16个不同频道的源，也可根据题设数量上限排除其他未知频道，但仍须清除所有已发现源。只清除了10至15个源时，虽然已经达到数量下限，仍可能存在额外源，必须继续排查。这也是少于16个源的场景有时仍需花费较长时间完成末尾搜索的原因。

\nextpage
\subsection{问题三的离线结果}
表\ref{tab:q3}列出本文方法的15场固定离线输入，合计192个源，均获得实际清除成功响应并正常退出。常规、边界、密集与备用场景分别列出；边界、密集和备用场景的分布及执行方式不同，不把它们混合为一个随机总体的性能估计。

\begin{table}[H]\centering\caption{问题3的15场离线固定输入结果}\label{tab:q3}
\small\setlength{\tabcolsep}{7pt}
\begin{tabular}{crrrrr}\toprule
案例 & 源数 & 总时间/s & 平均时间/(s/源) & $M$ & $K_f$\\\midrule
M1 & 16 & 2940.55 & 183.78 & 90 & 0\\
M2 & 10 & 3072.89 & 307.29 & 130 & 0\\
V10 & 10 & 2785.99 & 278.60 & 113 & 0\\
V12 & 12 & 3084.86 & 257.07 & 111 & 0\\
V14 & 14 & 3243.22 & 231.66 & 115 & 0\\
V16 & 16 & 3318.80 & 207.42 & 115 & 3\\
F10 & 10 & 2948.40 & 294.84 & 127 & 2\\
F12 & 12 & 3160.66 & 263.39 & 111 & 1\\
F14 & 14 & 2921.20 & 208.66 & 114 & 4\\
F16 & 16 & 3124.71 & 195.29 & 96 & 1\\
B10 & 10 & 4402.57 & 440.26 & 144 & 2\\
B16 & 16 & 4587.68 & 286.73 & 125 & 5\\
C10 & 10 & 2085.60 & 208.56 & 91 & 0\\
C16 & 16 & 437.67 & 27.35 & 33 & 0\\
R10 & 10 & 15348.44 & 1534.84 & 660 & 249\\
\bottomrule\end{tabular}
\end{table}


表中S1至S10为常规输入，B为边界输入，C为密集输入，R为备用过程；B、C、R后的数字为源个数。各案例的频道、位置、半径、方向和噪声设置见支撑材料中的合成案例数据。$M$表示测向次数，$K_f$表示光学失败次数，均由协议记录统计。

十场S类常规输入的平均总时间为3060.13秒。其中移动时间平均占比约76.01\%，说明单纯减少一次5秒测向通常不足以解释主要提速；任务次序、共享位置和清除出口对总费用更敏感。16源密集案例耗时437.67秒，比其他案例明显小，是因为源集中且达到数量上界后可以结束未知频道搜索，不表示源越多一定越容易。

强制备用案例耗时15348.44秒，显著高于主动调度。其作用是验证在前瞻和主动优化停止后仍能接续完成，不是估计正常调度的典型平均值。真实演练、正式测试以及现实运行窗口的结果另列，不用本表代替官方测试成绩。

\nextpage
\section{问题四：按运行顺序组织搜索、定位与清除}
问题四中，机器狗既不知道干扰源在哪里，也不知道各源向哪个方向发射。一次“没有信号”可能表示距离过远，也可能表示机器狗站在发射背面。因此，算法始终保留所有尚未被实际观测排除的位置，并围绕一个循环运行：\textbf{安排下一项任务，取得真实观测，更新源的位置范围，清除当前源，再安排剩余任务。}下面按这一执行顺序说明各步原理。

\subsection{运行开始：建立频道记录与搜索点布局}
机器狗从原点出发，为20个频道分别建立记录。开始时，各频道均为“未知”；首次收到信号后改为“已发现”，实际清除成功后改为“已清除”，能够确定不存在源时才改为“已排除”。不同频道的记录相互独立。在一个频道上没有收到信号，不能代替对其他频道的检查。

为使任何方向发射的源都有被发现的机会，程序准备21个搜索点：原点、半径998米的8个内环点，以及半径1868米的12个外环点。令$u(\theta)=(\cos\theta,\sin\theta)$，点集为
\begin{equation}\label{eq:q4stations}
\begin{split}
Q_4={}&\{(0,0)\}\cup\{998u(k\pi/4):k=0,\ldots,7\}\\
&\cup\{1868u(k\pi/6+\pi/12):k=0,\ldots,11\}.
\end{split}
\end{equation}
外环略超出1800米源域，是为了从外侧观察靠近边界、向外发射的源。内外环错开角度，则是为了让同一片区域能够从不同侧面被观察。

这些点的作用可以用“在近处围住源”解释。对任意可能的源位置$x$，取距它不超过1000米的一组测点$Q_x$，要求
\begin{equation}\label{eq:dircover}
\norm{q-x}\le1000\quad(q\in Q_x),\qquad x\in\conv(Q_x).
\end{equation}
后一个条件表示$x$落在这些点围成的凸多边形中。此时，无论源朝哪个半平面发射，测点都不可能全部落在它的背面，至少有一个位于可见侧；而1000米不超过任何合法源的接收半径，所以该点能够收到信号。只保证“附近有一个点”不够，这个点恰好在背面时仍会漏检。

程序用小方盒逐步细分源域，检查每个盒是否同时满足距离和包围条件，以此判断一组搜索点是否足够。未通过的盒继续细分；计算额度用完仍未通过时，保留原搜索任务。该检查随后也用于判断能否删去或替换某个尚未访问的站点。

\subsection{每轮开始：安排搜索站与已知源的处理顺序}
每轮首先整理两类尚未完成的任务：一类是到搜索站检查未知频道，另一类是定位并清除已经发现的源。把两类任务一起排序，可以在顺路发现源时及时处理它，减少“先走完全部搜索点，再折返清除”的往返距离。

每个任务用三个量表示：首先到达的入口、完成该任务的大致耗时、完成后的预计位置。搜索站的入口和出口都是站点本身，站内费用按待测频道数计算；源任务若已能直接清除，入口就是清除点，否则提供若干测向入口，并暂用源域中心估计出口。程序比较不同任务顺序及入口组合，选择移动与操作总费用较小的方案。

排序只确定当前先做什么。搜索任务执行完后，程序根据新发现的源重新排序；源任务一旦开始，通常持续处理到该源清除，再重排其余任务。这样既利用后续行程选择入口，又能及时响应新观测带来的变化。

\subsection{到达搜索站：逐频道检测并建立源的位置范围}
选中搜索站后，机器狗到达该点，对仍未知且尚未在此点检查过的频道逐一测量；若当前已调谐的频道仍需检查，先测该频道以减少切频。每次只根据对应频道的实际返回更新记录，具体有以下三种情况。

\textbf{收到方向值。} 这表明该频道确有干扰源，但测向误差使源不一定在示向直线上。算法以检测点为起点，保留示向两侧完整误差范围内的前向区域，并结合最大接收距离和源域范围，得到包含真实源的保守多边形$\widehat P_c$。程序用$1.005^\circ$包住题设$1^\circ$误差，给数值计算留出余量。后续得到新方向时，将新区域与旧区域求交，而不是直接把两条示向线的交点当作源位置。

\textbf{返回近距。} 这表示源已在当前点5米内。程序用边长10米的方盒保守表示这一信息。该范围远小于20米清除半径，因此随后的任务选择可直接转向光学清除。

\textbf{没有信号。} 该频道暂时仍为未知，只保存本次检测位置。单个阴性不能证明附近无源，因为源可能朝另一侧发射。只有该频道积累的实际阴性点已经足以覆盖所有位置与发射方向，才能把它改为“已排除”。计划以后去的站点不算已经取得的观测。

首次发现源时，算法还把其位置区域划分成若干小块，并给每块安排一个光学清除点，保证该块全部位于相应点的20米范围内。这些小块表示尚需排查的范围：以后每取得一条有效观测，就裁剪小块；即使测向难以继续，也可以沿这些清除点逐块排查。

\subsection{搜索结束：判断是否值得共享一次补测}
如果同一搜索站至少新发现三个源，且它们各自只有一次方向观测，那么一次短距离移动可能同时改善多个源的定位。程序在该站周围350米圆上，每隔$30^\circ$取一个候选观察点，比较哪些点值得绕行。

补测有用需要同时满足两个条件：\emph{到了那里有机会收到信号，收到的方向与原方向有足够夹角。}若两次观察方向几乎平行，即使收到信号，交会区域也可能仍然很长。设原区域沿测向轴的宽度为$w$，候选点到区域代表位置的距离为$d$，两次观察方向夹角为$\phi$，则补测后的宽度粗略估为
\begin{equation}
\widehat w'=\min\left\{w,\frac{2\varepsilon d}{\max(0.01,|\sin\phi|)}\right\},
\qquad \varepsilon=1.005\pi/180.
\end{equation}
分母体现视差的作用：夹角越小，横向测角误差越难转化为沿轴距离约束。

接收机会由该频道的全部正、负观测估计。程序用剩余小块的代表位置及面积近似未知源位置；对每个位置，保留1000至1500米内能够解释历史观测的接收半径。对定向源，还保留能朝向所有阳性点、并在距离足够时背向阴性点的发射方向。若一个阴性点本来就在假设接收半径之外，它不限制朝向。由此得到候选点的接收概率估计$\widehat p_c(q)$。

将宽度减少量的一半近似折算成可节省的移动距离，再按5米每秒换算，并扣除测向与可能切频的6秒，得到
\begin{equation}\label{eq:sharedscore}
\widehat G_c(q)=\widehat p_c(q)\frac{w-\widehat w'}{10}-6.
\end{equation}
仅为正收益频道安排补测。所有正收益之和还须超过绕行时间：从搜索站$p$绕到$q$后再去参考目标$t$，比直接去$t$多用$(350+\norm{q-t}-\norm{p-t})/5$秒。这里$t$取最近的已知源域中心，用于估计后续去向。

上述概率采用全向与定向各半、半径与朝向均匀的工作假设，只用于比较行动。实际补测若没有信号，仍记为阴性，不按预测结果缩小位置区域。

\subsection{开始处理一个源：先判断是否已经能够清除}
任务排序选中源$c$后，先问“现在能否保证进入它的20米范围”，而不要求先求出一个高精度坐标。若存在位置$q$满足
\begin{equation}
\max_{x\in\widehat P_c}\norm{q-x}\le19.8,
\end{equation}
则不论源在当前区域的哪里，从$q$都能完成清除。程序检查多边形顶点即可判断整个区域是否位于该圆内，并从可行清除位置中选择移动代价较小者。19.8米比题设20米略小，用于留出余量；若已有下一任务的入口，还同时考虑清除后去往该入口的距离。

如果暂时没有这样的清除点，就先执行任务排序给出的测向入口。此后每次取得结果，都更新区域，再检查能否直接清除；只有仍不能清除时，才继续选择下一次测向或光学排查。这里不断缩小的是“源可能在哪里”的范围，而不是反复修正一个被当作真实位置的估计点。

\subsection{需要继续测向：在原示向两侧选择检测点}
继续测向时，以最近一次确实收到方向值的位置$o$为参考，令$u$为示向单位向量，$v$为其垂直方向。首次定位阶段，在沿轴前进的不同距离上，分别向两侧偏移，形成一对检测点
\begin{equation}
q_{\pm}=o+\bigl[5+(L_c-5)\lambda\bigr]u\pm(L_c/50)v,
\qquad\lambda\in\{0.38,0.50,0.62\},
\end{equation}
其中$L_c$是当前从$o$到源的距离上界。沿轴前进是为了接近源，两侧偏移是为了获得新的交会角度，并为一侧处于发射背面的情况准备另一侧的检测位置。

取得至少两次方向观测后，算法把当前区域投影到最新示向轴上，得到剩余区间$[\ell,h]$。新的前进位置改用区间中点$(\ell+h)/2$，侧向距离取$\max\{5,(h-\ell)/20\}$米。这样会随已经缩小的区域调整步长，不再从首次观测的完整距离范围重新搜索。

一个源刚被选中时，测点及先后次序由整体任务排序给出。源内继续定位时，程序比较候选点对，并估算“测量以后还剩多少排查工作”：把可能返回的方向分成六个角区间，对各区间对应的剩余区域估算光学排查费用，同时计入近距和双阴性分支。选择时兼顾到点移动与后续工作量，避免只因某点最近就前往。

该估算受计算时间限制。区域仍包含较多小块时，程序也允许继续尝试测向，以避免过早执行很长的光学序列；但测向轮数和额外试探次数均有限。若拟采用的一对点已经实际得到过双阴性，则改用光学排查。同位置的电磁环境固定，原样重复这一对测量不会产生新的信息。

\subsection{按检测返回更新区域：第一点无信号才试第二点}
对选定的一对点，先到第一点测量。若得到方向值，就与已有区域求交，并以该次观测作为新的参考；若得到近距结果，就按近距范围更新。两种情况都已经取得有效定位信息，本轮不再机械地走到第二点。

只有第一点无信号时，才前往另一侧的第二点。第二点若收到信号，就正常更新；两点都无信号时，再结合原来确实收到信号的参考点$o$，判断能否排除一部分候选位置。

\subsection{双阴性之后：只删除无法解释真实观测的位置}
两个阴性点$a,b$也不能无条件删去它们附近的区域。算法需要先确认：对于准备删除的那一片区域，任意候选源位置$x$到$a,b$都不超过1000米。满足这一条件后，“距离太远”就不能解释两次阴性；如果源是定向的，只能是$a,b$都在其发射背面。

接着检查相对方向。如果从$x$看去，阳性点$o$的方向落在$a,b$两个方向张成的非负锥内，即
\begin{equation}
o-x=\alpha(a-x)+\beta(b-x),\qquad \alpha,\beta\ge0,
\end{equation}
那么任何同时背向$a,b$的发射半平面，也不可能朝向$o$。这与$o$曾经实际收到信号矛盾，所以这样的$x$可以排除。若源为全向，两个1000米内测点无信号本身就已经矛盾。

\fig{directional_geometry}{0.99}{定向观测的两种几何关系。左图中，附近测点围住源，使任意发射半平面内至少存在一个测点。右图展示双阴性后的候选排除区域：只有整片区域同时满足两点距离约束和相对方向约束，才能删除。图为原理示意。}

实现时，程序先取$a,b$连线远离$o$一侧的候选部分，再对其各顶点检查上述距离与方向条件；整片通过才进行裁剪。条件不满足时，就保留位置区域，转而选择别的测向或光学行动。例如，位于$o,a,b$三角形内部的源完全可能朝向$o$、背向$a,b$，因此不能因为两次阴性就删去三角形内部。

这一步的作用是利用定向阴性中确实包含的信息，而不是强行让每次阴性都产生区域缩小。即使区域没有缩小，本次阴性位置仍保留在记录中，后续估计哪个观察点更容易收到信号时会用到它。

\subsection{下一轮源内决策：考虑用短光学序列代替继续测向}
新的方向观测或有效的阴性收缩，都可能使位置区域变成窄带。若此时一个清除点仍不能覆盖整个区域，继续绕到可见侧测向可能不划算。光学探测不受源的发射方向限制，因此后续源内决策还会尝试用2至8个清除点覆盖整个区域，组成一条短光学序列。

构造时分别沿最新示向轴及其垂直方向摆放点。窄带越宽，一个20米圆能沿长轴覆盖的长度就越短，所需点数也越多。程序先按略小于19.8米的圆半径估算点间距，再逐一检查完整区域是否被覆盖；无法形成合格短序列时，继续使用测向或原有的小块排查方案。

\subsection{执行光学序列：先确认完整覆盖，再比较先后次序}
光学候选点$Q=(q_1,\ldots,q_k)$必须满足
\begin{equation}\label{eq:optcover}
\widehat P_c\subseteq\bigcup_{i=1}^{k}B(q_i,19.8).
\end{equation}
程序不是只检查几个代表位置，而是把整个多边形按“离哪个清除点最近”分成Voronoi小块。对属于$q_i$的小块，只要其所有顶点都在$q_i$的19.8米圆内，整个小块就在圆内；各块合起来覆盖原区域，因而不会遗漏块与块之间的位置。

\fig{optical_chain}{0.62}{一次离线Q4决策中的完整光学序列。黑色边界为当时的保守位置区域，虚线圆半径为19.8米，填色区表示分配给不同清除点的Voronoi小块。数字为候选访问次序。图只使用当时已有观测，不使用源的真实坐标。}

覆盖通过后，再决定从哪一点开始。因为成功后即可停止，所以完整走完的路线最短，不一定意味着通常最快找到源。程序比较沿轴正序、逆序，以及优先访问“单位移动与探测费用能覆盖更多剩余面积”的顺序。用小块代表位置和面积估计首次在第$i$点成功的概率$\widehat\pi_i$，相应费用为
\begin{equation}\label{eq:optmean}
\widehat T_o(Q)=\sum_{i=1}^k\widehat\pi_i\left[\frac{\norm{p-q_1}+\sum_{j=2}^{i}\norm{q_{j-1}-q_j}}5+3(i-1)+5\right].
\end{equation}
括号中的三项依次是走到首个成功点前的移动时间、前面$i-1$次失败探测各3秒、最后一次探测与清除共5秒。若已确定下一任务入口，还加上从成功点前往该入口的移动时间。

程序将较好的光学序列与继续测向的费用估计比较。后者乐观地假设下一次测向能提供有效信息，再估计接近源并清除的费用。光学估计更小时，且整条序列的行动预算足够，就执行光学序列。面积权重用于安排先后次序；保证能够覆盖所有可能位置的依据仍是式\eqref{eq:optcover}。

执行时逐点调用真实清除接口。某点失败，只删除被该点20米圆完整包含的旧小块，其他小块保留；不能把一次失败解释为整个频道无源。一旦返回清除成功，立即取消该源后续所有动作，将频道改为“已清除”。因此，计划准备的是完整序列，实际执行到第一次成功就结束。

\subsection{完成当前任务：利用当前位置补测，再重排剩余任务}
完成一个源的清除后，机器狗不立即离开。程序先检查当前位置是否适合补测其他已发现源；搜索站的一批共享补测结束后，也执行相同检查。已经走到的位置无需额外移动，若观察角度合适，一次切频测量就可能减少下一源的定位工作。

原地补测仍使用式\eqref{eq:sharedscore}比较接收机会、预计区域缩小和6秒操作费用。问题四还要求当前点的999米圆包含该源的整个保守区域，使“距离足够近”成立；同时要求视差$|\sin\phi|\ge0.25$，且此点未对该源测过。优先处理净收益较大的频道，避免为了几乎平行的新方向反复切频。已有至少三次阳性观测的源不再进入这一原地补测候选集。

补测完成后，回到整场调度。刚清除的源不再占用任务，其他源域也可能因顺路观测而缩小。程序据此重新安排剩余搜索站与源任务。若当前清除位置能够替代某个待搜索站，会对“所有未知频道共同测过的历史点，加上新的待访问点”重新检查覆盖；只有仍足够发现所有合法源时才调整路线。替换成功只改变未来计划，对未知频道的排除仍需等实际检测结果。

\subsection{额度或时间不足：接续执行有限的备用排查}
在每项额外行动之前，算法检查执行该行动后是否仍有足够时间完成剩余排查。额外试探额度、现实时间或虚拟时间余量不足时，就停止继续比较优化方案，转入预先保留的备用过程。额度贯穿整场累计，不因处理下一个源而重新开始。

备用过程先清除已经发现的源，按剩余小块对应的清除点逐块排查。随后沿间距600米的$7\times7$网格，逐点检测尚未解决的未知频道。每个网格方格的四角都在格内任意位置的$600\sqrt2<1000$米范围内，并且四角围住该位置，所以同样能处理方向未知的源。新发现一个源，就先完成它的光学排查，再继续网格搜索。

这一过程能够接续前面的工作，是因为真实观测只删去已经排除的范围，剩余小块始终覆盖还可能存在源的位置。阴性检查逐步减少未查网格点，光学失败逐步减少待查小块，成功则结束一个源任务；有限的额外试探不会无限重复。在题设物理规则、有效通信和足够执行窗口下，这为主动选择难以继续改善的情形保留了完成途径。

\subsection{整场结束：同时满足清除完成与未知频道排除}
程序每轮都检查式\eqref{eq:stop}。结束需要同时满足：已经发现的源全部取得实际清除成功响应，所有其他频道均已排除，且实际清除数处于题设的10至16个范围内。

未知频道有两种排除依据。一种是该频道在足以发现任意合法源的实际检测点上始终无信号；另一种是已发现16个不同频道的源，达到了题设数量上限。达到16个时可以停止寻找更多源，但仍需把已发现源全部清除。若只清除了10至15个，则还要继续排查剩余未知频道，不能仅凭已经达到最少源数结束。

\nextpage
\subsection{问题四的离线结果与波动解释}
当前15场Q4离线输入合计198个源，均实际清除。表\ref{tab:q4}报告总时间、平均每源时间及光学失败次数，含全定向、混合、边界、密集和强制备用输入。所有输入由独立模拟器持有，求解器只访问协议响应。

\begin{table}[H]\centering\caption{问题4的15场离线固定输入结果}\label{tab:q4}
\small\setlength{\tabcolsep}{7pt}
\begin{tabular}{ccrrr}\toprule
案例 & 全向/定向 & 总时间/s & 平均时间/(s/源) & $K_f$\\\midrule
M1 & 5/11 & 4551.45 & 284.47 & 1\\
M2 & 2/8 & 6229.88 & 622.99 & 2\\
V10 & 0/10 & 6365.14 & 636.51 & 0\\
V12 & 0/12 & 6534.05 & 544.50 & 18\\
V14 & 0/14 & 6399.30 & 457.09 & 1\\
V16 & 0/16 & 6680.12 & 417.51 & 15\\
F10 & 4/6 & 6381.56 & 638.16 & 2\\
F12 & 4/8 & 5939.68 & 494.97 & 1\\
F14 & 5/9 & 6359.10 & 454.22 & 3\\
F16 & 6/10 & 7176.66 & 448.54 & 20\\
B10 & 0/10 & 8060.39 & 806.04 & 16\\
B16 & 0/16 & 9206.85 & 575.43 & 28\\
C10 & 0/10 & 7148.16 & 714.82 & 8\\
C16 & 0/16 & 4290.26 & 268.14 & 37\\
R16 & 0/16 & 17850.63 & 1115.66 & 893\\
\bottomrule\end{tabular}
\end{table}


十场S类常规输入平均总时间为6261.69秒，移动时间平均占比约74.01\%。与问题三相比，定向约束使接近源以后仍可能无法测向，进入新视角和排除未知频道的成本随之增加。

边界16源案例耗时9206.85秒，密集16源案例耗时4290.26秒。相同源数对应不同用时，说明还需结合源的位置、发射方向和具体观测过程解释效率。达到16源允许停止剩余频道检查，但这一好处未必抵消此前定位与移动的费用。

光学序列的先后安排依赖面积权重和后续去向估计，实际源若落在较晚访问的小块内，清除用时仍会增加。后文保留同输入消融与正式测试表，用于补充各步骤的实际费用和作用。

\nextpage
\section{行动成本与时间分配}
\subsection{移动与设备操作的统一计量}
每个动作的总成本由移动、测向、切频及光学操作组成。移动速度为5米每秒，距离$d$对应$d/5$秒；测向5秒，切换频道另加1秒，光学成功5秒，失败3秒。程序使用模拟器返回的累计虚拟时间更新运行状态，现实时间另用于限制规划和接口等待。

\fig{task_costs}{0.97}{行动费用的组成。左图给出移动距离与时间的线性关系；右图给出四种设备操作的单次费用。纵轴均为秒，可直接用于同一任务序列的费用累计。}

移动100米需要20秒，已超过一次测向或一次清除的设备费用。因此，减少跨源往返、合并共享观察点、选择合适的服务出口，是控制总费用的主要途径。设备操作仍不可忽略：未知频道较多时，同站逐频道排查会累计测向及切频成本。

\subsection{在线计算与备用过程的衔接}
射频候选的默认规划窗口为0.2秒，可选动作初始额度为640次。额度用于限制没有减少剩余任务的试探。当候选计算结束时，立即执行已获得的可用方案；不等待所有可能的后续动作都被枚举。

行动选择还需计入剩余任务的完成成本。程序按尚未排查的网格点、已发现源的剩余清除位置及可能未发现的源数估计备用工作量。若主动动作所需费用与备用过程合计超过剩余预算，则转入备用过程。现实窗口紧张时也采用同一转移，避免把剩余时间耗在规划上。

备用过程先处理已发现源，再沿固定网格逐频道搜索。每次成功清除后释放对应任务；实际没有收到信号时，只更新该频道的记录。初始额度、现实时间和任务计数均随整场运行累计，不因重新规划而重置。

时间分配同时影响效率与完成能力。更长的规划窗口可能改善局部选择，也会占用现实时间；更多试探可能减少移动，也可能产生无效测量。因此参数影响需要在相同输入下单独测量，相关未测结果保留在后面的参数实验表中。

\nextpage
\section{计算结果与成本分析}
\subsection{整体完成情况}
30场合成案例包含常规、边界、密集和备用场景，合计390个干扰源，均获得实际清除成功响应并退出。问题三和问题四的完整分场结果分别见表\ref{tab:q3}和表\ref{tab:q4}。这些是合成环境中的计算结果，正式测试结果另列。

\fig{performance}{0.99}{常规案例的用时与费用构成。左图展示各题10场案例的源个数和总时间；右图将平均总时间分解为移动、测向、切频及光学清除费用。同源数的案例具有不同位置和接收条件，散点不表示同一随机量的重复测量。}

\begin{table}[H]\centering\caption{常规案例的分题汇总}\small
\begin{tabular}{crrr}\toprule
问题 & 案例数 & 平均总时间/s & 平均移动时间占比\\\midrule
三 & 10 & 3060.13 & 76.01\%\\
四 & 10 & 6261.69 & 74.01\%\\\bottomrule
\end{tabular}\end{table}

两题的移动费用均占总时间的主要部分。问题四中，定位点还必须位于可接收的一侧，几何距离较短并不意味着测向有效。一次阴性结果可能来自接收半径，也可能来自方向限制；无法立即排除大片候选区域时，需要增加视角或改用光学清除。

源个数并不能单独决定总时间。达到16个已发现源时，可以利用数量上限排除剩余未知频道；少于16个时仍需完成这些频道的排查。位置集中、可接收方向、源间访问距离与最终未知频道排除共同决定时间变化。

\nextpage
\subsection{参数敏感性与模块作用}
边界与密集输入分别考察接近源域边缘及局部集中分布的情况。两类输入使用固定的极端正负一度误差，同地点重复测量不会产生可平均的新噪声样本。其结果应与常规输入分开分析。

评价参数作用时，需要固定源位置、半径、方向和噪声，仅改变一个参数。下表保留尚未完成的结果，空栏不表示零变化。

\begin{table}[H]\centering\caption{待完成的参数敏感性实验}\small
\begin{tabularx}{\linewidth}{YYcc}\toprule
参数 & 比较水平 & Q3平均时间/s & Q4平均时间/s\\\midrule
共享观察半径 & 250、350、450 m & \blank & \blank\\
射频前瞻预算 & 0.05、0.20、0.50 s & \blank & \blank\\
初始可选额度 & 320、640、960次 & \blank & \blank\\
安全清除半径 & 19.5、19.8、19.9 m & \blank & \blank\\\bottomrule
\end{tabularx}\end{table}

每组都须同时报告全清比例、最差案例时间、现实运行时间和备用触发次数。若缩小安全余量带来时间下降，还必须重新检查连续覆盖与数值边界，不能只比较平均时间。若前瞻预算增加而虚拟时间不变，应检查是否原候选已足够、是否决策主要由备用条件控制，或额外计算没有改变所选动作。

\begin{table}[H]\centering\caption{待完成的同输入消融实验}\small
\begin{tabularx}{\linewidth}{Yccc}\toprule
对照设置 & 全清比例 & 相对完整策略$\Delta T$/s & 现实时间/s\\\midrule
关闭原地补测 & \blank & \blank & \blank\\
关闭共享第二测点 & \blank & \blank & \blank\\
关闭任务重排，仅固定站序 & \blank & \blank & \blank\\
关闭2--8点光学链，保留完整备用 & \blank & \blank & \blank\\
Q4接收评分不使用历史阴性 & \blank & \blank & \blank\\\bottomrule
\end{tabularx}\end{table}

预先约定$\Delta T=T_{\text{对照}}-T_{\text{完整}}$，正值表示该模块在对应输入中有收益。两种方法必须使用相同案例，并保留失败和超时；模块间有相互作用，单项收益不能简单相加。上述分析方法已经确定，具体效应大小留待实际运行后填写。

\nextpage
\section{官方演练汇总与正式测试安排}
\subsection{已经提供的演练记录}
参赛队保存了两题各五份演练日志及按运行顺序排列的截图。当前可核对日志的案例编码、题号及时间等公开封装信息，正文为加密数据；尚未从中解析逐步动作与真实源布局。因此表\ref{tab:practice}仅记载截图报告的源个数和整数总时间，不声称已独立解密验证全清比例。

\begin{table}[H]\centering\caption{十次演练截图的成绩汇总}\label{tab:practice}
\small\setlength{\tabcolsep}{4pt}\begin{tabular}{cccrr}\toprule
题号/次序 & 案例编码 & 报告全向/定向 & 报告个数 & 总时间（约）/s\\\midrule
3/1 & 4GC8-BUAN-JTNR-732X & 16/0 & 16 & 3233\\
3/2 & ZAJ4-Y9H4-X94F-84RF & 11/0 & 11 & 2595\\
3/3 & PXX5-MVVF-4U5N-RJ29 & 15/0 & 15 & 3123\\
3/4 & YVXK-H49D-HS89-HBGH & 12/0 & 12 & 2826\\
3/5 & VD4K-3FT8-EMN3-KHMN & 11/0 & 11 & 3189\\
4/1 & WZM3-U3MR-YW22-QHNK & 0/14 & 14 & 7570\\
4/2 & ZG4N-CQJP-AZ7F-6M34 & 14/2 & 16 & 4559\\
4/3 & 2AA5-BWCU-M8MH-M9A2 & 3/8 & 11 & 6439\\
4/4 & C2MT-ZV9Z-3G7M-YE4G & 2/9 & 11 & 7191\\
4/5 & 6QJH-B56T-UGBD-NJZ3 & 9/2 & 11 & 6249\\
\bottomrule\end{tabular}
\end{table}


\fig{practice}{0.82}{十次演练截图的描述性比较，点旁数字为运行顺序。总时间为截图取整值；Q4横轴为报告的定向源数。不同点来自不同案例，未控制位置、半径、方向和噪声，因此未拟合因果趋势线。}

问题四时间为4559至7570秒，跨度3011秒，大于问题三的638秒。定向较多可能增加绕行与无效测向，但源位置及数量也不同。全向9个、定向2个的案例仍耗时6249秒，说明仅用定向比例无法解释全部波动。需用同次运行的明文动作日志分解移动、测向、失败光学与末尾排除成本，再做同输入消融。

每次演练的实际清除比例和现实运行时间目前分别预留为\blank、\blank。只有补齐对应原始结果后，才能按题设定义计算平均定位清除时间，不将截图报告源数未经核对地当作清除成功数。

\nextpage
\subsection{问题三正式测试结果}
题目要求在充分演练后进行三次正式测试，并保留模拟器导出的原始日志\cite{problem,manual}。已取得的正式成绩列于下表，尚未取得的字段保留空栏。测试完成后据实填写相应字段，不用其他演练或离线案例的数值占位。

\begin{table}[H]\centering\caption{问题三正式测试结果（空栏待实测）}\small
\begin{tabularx}{\linewidth}{Yccc}\toprule
测试案例编码 & 清除源个数 & 平均定位清除时间/s & 程序运行时间/s\\\midrule
\blank & \blank & \blank & \blank\\[10pt]
\blank & \blank & \blank & \blank\\[10pt]
\blank & \blank & \blank & \blank\\[10pt]\bottomrule

\end{tabularx}\end{table}

\subsection{问题四正式测试结果}
\begin{table}[H]\centering\caption{问题四正式测试结果（空栏待实测）}\small
\begin{tabularx}{\linewidth}{Yccc}\toprule
测试案例编码 & 清除源个数 & 平均定位清除时间/s & 程序运行时间/s\\\midrule
\input{\PaperDir/formal_q4.tex}
\end{tabularx}\end{table}

\subsection{取得数据后的计算与分析}
对每次正式测试，先匹配案例编码、题号、策略标识和原始日志，核对实际清除成功数$K_s$及总虚拟时间$T_v$；再按$\overline T=T_v/K_s$计算平均定位清除时间，现实运行时间按模拟器给出的记录填写。只有存在实际源总数$N$的可信结果时，才另报$\eta=K_s/N$，不可用程序自身的完成声明替代外部真值。

三次结果的分析同时报告均值、最大值和取值范围。样本数仅为3，不能据此稳定估计极端尾部概率；若一场明显较慢，应保留该场，检查是否由移动、定向阴性、光学失败或未知频道排除构成主要费用。若发生超时或未清完，则报告实际完成数、结束原因和已消耗时间，不只对成功场计算平均表现。

六份正式日志的原始文件名分别待填入支撑材料清单，不改名、不补造文件。材料整理和论文编译过程均不启动正式模块。本节保留完整的数据取得、指标计算和结果分析框架，待实测数据到位后即可补写具体比较结论。

\nextpage
\section{模型评价与结论}
\subsection{主要优点}
\textbf{证据与费用分层。} 覆盖证书、实际成功响应及下降秩决定完成性，概率与路径代理决定动作次序。该结构允许近似评分存在误差，并以实际观测和清除结果更新任务状态。

\textbf{统一考虑定位与移动。} 第二点设计利用首次接收约束，多源排序考虑服务入口与出口，狭窄源域可直接转为完整光学链。这些处理把几何精度转化为实际任务费用，而不是无限追求更精确的坐标估计。

\textbf{对定向不可见作显式处理。} 21站布局通过连续凸包覆盖核查，单个阴性不被误当成距离证据，双阴性收缩要求真实观测前提。完整备用过程使主动策略在计算预算不足时仍有可接续路径。

\subsection{局限与可改进方向}
全清结论以题设物理条件、正确数值包络及足够执行窗口为前提；当前入口对通信异常直接停止，没有自动重试。接收评分、源出口和光学首次成功质量均为近似，尚缺独立校准。30场固定输入检验了实现行为，不能穷尽所有连续位置、半径和朝向；用于开发的输入也不能代替未见场景上的统计保证。

角度优化针对给定移动预算下的解析直径上界；整场调度仍使用有限候选和近似费用。故本文不宣称已逼近连续原问题的全局理论最优。后续可在保持证据账本不变的前提下，先做同输入模块消融，再提高服务费用与接收模型的准确度，并研究包含移动和信息获取的更强下界。

\subsection{结论}
问题一建立了前向楔形交直径算法，并否定同直径圆必然覆盖的命题。问题二得到保证接收的三圆盘候选域，并在给定移动预算下通过角区间二分确定第二测点。问题三、四分别采用七站与21站的连续发现依据，将频道证据、联合任务排序、共享观测和光学清除连接为在线策略。

30场合成案例中390个源全部清除。常规输入中，问题三、四平均总时间分别为3060.13秒和6261.69秒；不同位置与方向组合产生明显波动。正式测试与尚未完成的敏感性、消融结果均保留待填字段。当前结论由已取得的计算和观测结果支持，新增实测将用于补充效率与可靠性评价。

\section*{AI工具使用声明}
本参赛队在竞赛过程中使用了AI工具，主要用于算法讨论、代码编写与调试、资料整理、数值核查、图表生成和论文初稿撰写，详细使用情况见支撑材料。本稿的人工修改与逐项核验情况尚需参赛队据实补充，不将自动化测试通过等同于人工审查完成\cite{airules}。

\nextpage
\begin{thebibliography}{9}
\bibitem{problem} 全国大学生数学建模竞赛组委会. 2026年高教社杯全国大学生数学建模竞赛B题：无线电干扰源的快速自动定位与清除[Z]. 2026.
\bibitem{manual} 全国大学生数学建模竞赛组委会. 2026年B题附件1：模拟器测试说明[Z]. 2026.
\bibitem{interface} 全国大学生数学建模竞赛组委会. 2026年B题附件2：模拟器通信接口说明及编程指南[Z]. 2026.
\bibitem{toussaint} TOUSSAINT G T. Solving geometric problems with the rotating calipers[C]//IEEE MELECON'83. Athens, 1983: A10.02/1--4. \url{https://www-cgrl.cs.mcgill.ca/~godfried/research/calipers.html}.
\bibitem{boyd} BOYD S, VANDENBERGHE L. Convex Optimization[M]. Cambridge: Cambridge University Press, 2004. \url{https://stanford.edu/~boyd/cvxbook/}.
\bibitem{kuhn} KUHN H W. The Hungarian method for the assignment problem[J]. Naval Research Logistics Quarterly, 1955, 2(1--2): 83--97. DOI:10.1002/nav.3800020109.
\bibitem{airules} 全国大学生数学建模竞赛组委会. 全国大学生数学建模竞赛人工智能工具使用规定（2026年试行）[EB/OL]. (2026-08-03)[2026-09-13]. \url{https://www.mcm.edu.cn/html_cn/node/fef94648f2836ab6cc81586f4c38512b.html}.
\end{thebibliography}

\label{main_end}
\clearpage
\begin{center}{\Large\heiti\bfseries 附\quad 录}\end{center}
\section*{附录1：支撑材料文件列表}
表中路径相对于支撑材料压缩包根目录，以“/”结尾的项为目录。正式日志待实测补入，保留原始文件名及内容；共用程序只在对应附录列出一次。
\begingroup\small\linespread{1.0}\selectfont
\setlength{\tabcolsep}{5pt}\renewcommand{\arraystretch}{1.18}
\begin{longtable}{@{}>{\raggedright\arraybackslash}p{.49\linewidth}>{\raggedright\arraybackslash}p{\dimexpr.51\linewidth-10pt\relax}@{}}
\caption{支撑材料文件列表}\label{tab:materials}\\
\toprule[1pt]\multicolumn{1}{c}{文件名} & \multicolumn{1}{c}{说明}\\\midrule\endfirsthead
\multicolumn{2}{c}{表\thetable\quad 支撑材料文件列表（续）}\\
\toprule[1pt]\multicolumn{1}{c}{文件名} & \multicolumn{1}{c}{说明}\\\midrule\endhead
\bottomrule[1pt]\endfoot
\texttt{\detokenize{AI工具使用详情.pdf}} & AI工具使用环节、用途与人工核验待填记录\\
\path{paper/results.json} & 已有离线计算结果及演练成绩汇总\\
\path{submission_inputs/practice_results.csv} & 问题三、四演练成绩；未取得的字段留空\\
\path{submission_inputs/formal_results.csv} & 问题三、四各三次正式成绩表，待实测填写\\
\path{paper/cases.json} & 30场离线合成案例输入\\
\path{paper/figure-data.json} & 光学清除图的可行域和清除点数据\\
\path{solve_questions.py} & 问题一、二的JSON输入输出入口\\
\path{src/jammer_solver/questions.py} & 问题一交会区域计算与问题二第二测点选择\\
\path{run_robot.py} & 问题三、四的模拟器运行入口\\
\path{start_robot.ps1} & Windows交互启动脚本\\
\path{src/jammer_solver/solver.py} & 多源搜索、定位、调度与清除算法\\
\path{src/jammer_solver/exact_geometry.py} & 可行域与光学覆盖所用几何运算\\
\path{src/jammer_solver/protocol.py} & 模拟器通信与运行记录\\
\path{src/jammer_solver/__init__.py} & Python包标记文件，内容为空\\
\path{paper/make_tables.py} & 生成离线结果和演练成绩表\\
\path{paper/make_formal_tables.py} & 从正式成绩CSV生成论文表格\\
\path{paper/make_figures.py} & 绘制论文图形\\
\path{paper/generate_appendix.py} & 生成材料文件列表及完整程序附录\\
\path{paper/build.py} & 编译论文与AI工具使用详情\\
\path{prepare_submission.py} & 打包支撑材料并列出待补事项\\
\path{README.md} & 材料目录及使用说明\\
\path{docs/running.md} & 安装、运行与模拟测试操作说明\\
\path{docs/algorithm.md} & 算法流程与模块说明\\
\path{paper/} & 论文源文件、表格、AI说明及编译说明\\
\path{paper/figures/} & 正文使用的8幅PDF图形\\
\path{submission_inputs/README.md} & 待补结果的字段和存放规则\\
\path{submission_inputs/ablation/} & 消融实验材料预留目录，待补\\
\path{submission_inputs/sensitivity/} & 敏感性实验材料预留目录，待补\\
\path{submission_inputs/practice_details/} & 演练清除数量与现实运行时间明细，待补\\
\path{submission_inputs/formal/q3/} & 问题3第1、2、3次正式日志待补\\
\path{submission_inputs/formal/q4/} & 问题4第1、2、3次正式日志待补\\
\end{longtable}
\endgroup
\clearpage
\section*{附录2：问题一、二代码}
\subsection*{\texttt{\detokenize{solve_questions.py}}}
问题一、二的JSON输入输出入口。
\begingroup\linespread{1.0}\selectfont\VerbatimInput[fontsize=\fontsize{8.5}{10.5}\selectfont,breaklines=true,breakanywhere=true,numbers=left,numbersep=4pt,xleftmargin=17pt,tabsize=4]{\RepoRoot/solve_questions.py}\endgroup
\subsection*{\texttt{\detokenize{src/jammer_solver/questions.py}}}
问题一交会区域计算与问题二第二测点选择。
\begingroup\linespread{1.0}\selectfont\VerbatimInput[fontsize=\fontsize{8.5}{10.5}\selectfont,breaklines=true,breakanywhere=true,numbers=left,numbersep=4pt,xleftmargin=17pt,tabsize=4]{\RepoRoot/src/jammer_solver/questions.py}\endgroup
\clearpage
\section*{附录3：问题三、四代码}
\subsection*{\texttt{\detokenize{run_robot.py}}}
问题三、四的模拟器运行入口。
\begingroup\linespread{1.0}\selectfont\VerbatimInput[fontsize=\fontsize{8.5}{10.5}\selectfont,breaklines=true,breakanywhere=true,numbers=left,numbersep=4pt,xleftmargin=17pt,tabsize=4]{\RepoRoot/run_robot.py}\endgroup
\subsection*{\texttt{\detokenize{start_robot.ps1}}}
Windows交互启动脚本。
\begingroup\linespread{1.0}\selectfont\VerbatimInput[fontsize=\fontsize{8.5}{10.5}\selectfont,breaklines=true,breakanywhere=true,numbers=left,numbersep=4pt,xleftmargin=17pt,tabsize=4]{\RepoRoot/start_robot.ps1}\endgroup
\subsection*{\texttt{\detokenize{src/jammer_solver/solver.py}}}
多源搜索、定位、调度与清除算法。
\begingroup\linespread{1.0}\selectfont\VerbatimInput[fontsize=\fontsize{8.5}{10.5}\selectfont,breaklines=true,breakanywhere=true,numbers=left,numbersep=4pt,xleftmargin=17pt,tabsize=4]{\RepoRoot/src/jammer_solver/solver.py}\endgroup
\subsection*{\texttt{\detokenize{src/jammer_solver/exact_geometry.py}}}
可行域与光学覆盖所用几何运算。
\begingroup\linespread{1.0}\selectfont\VerbatimInput[fontsize=\fontsize{8.5}{10.5}\selectfont,breaklines=true,breakanywhere=true,numbers=left,numbersep=4pt,xleftmargin=17pt,tabsize=4]{\RepoRoot/src/jammer_solver/exact_geometry.py}\endgroup
\subsection*{\texttt{\detokenize{src/jammer_solver/protocol.py}}}
模拟器通信与运行记录。
\begingroup\linespread{1.0}\selectfont\VerbatimInput[fontsize=\fontsize{8.5}{10.5}\selectfont,breaklines=true,breakanywhere=true,numbers=left,numbersep=4pt,xleftmargin=17pt,tabsize=4]{\RepoRoot/src/jammer_solver/protocol.py}\endgroup
\subsection*{\texttt{\detokenize{src/jammer_solver/__init__.py}}}
Python包标记文件，内容为空。
\clearpage
\section*{附录4：图表与材料生成代码}
\subsection*{\texttt{\detokenize{paper/make_tables.py}}}
生成离线结果和演练成绩表。
\begingroup\linespread{1.0}\selectfont\VerbatimInput[fontsize=\fontsize{8.5}{10.5}\selectfont,breaklines=true,breakanywhere=true,numbers=left,numbersep=4pt,xleftmargin=17pt,tabsize=4]{\RepoRoot/paper/make_tables.py}\endgroup
\subsection*{\texttt{\detokenize{paper/make_formal_tables.py}}}
从正式成绩CSV生成论文表格。
\begingroup\linespread{1.0}\selectfont\VerbatimInput[fontsize=\fontsize{8.5}{10.5}\selectfont,breaklines=true,breakanywhere=true,numbers=left,numbersep=4pt,xleftmargin=17pt,tabsize=4]{\RepoRoot/paper/make_formal_tables.py}\endgroup
\subsection*{\texttt{\detokenize{paper/make_figures.py}}}
绘制论文图形。
\begingroup\linespread{1.0}\selectfont\VerbatimInput[fontsize=\fontsize{8.5}{10.5}\selectfont,breaklines=true,breakanywhere=true,numbers=left,numbersep=4pt,xleftmargin=17pt,tabsize=4]{\RepoRoot/paper/make_figures.py}\endgroup
\subsection*{\texttt{\detokenize{paper/generate_appendix.py}}}
生成材料文件列表及完整程序附录。
\begingroup\linespread{1.0}\selectfont\VerbatimInput[fontsize=\fontsize{8.5}{10.5}\selectfont,breaklines=true,breakanywhere=true,numbers=left,numbersep=4pt,xleftmargin=17pt,tabsize=4]{\RepoRoot/paper/generate_appendix.py}\endgroup
\subsection*{\texttt{\detokenize{paper/build.py}}}
编译论文与AI工具使用详情。
\begingroup\linespread{1.0}\selectfont\VerbatimInput[fontsize=\fontsize{8.5}{10.5}\selectfont,breaklines=true,breakanywhere=true,numbers=left,numbersep=4pt,xleftmargin=17pt,tabsize=4]{\RepoRoot/paper/build.py}\endgroup
\subsection*{\texttt{\detokenize{prepare_submission.py}}}
打包支撑材料并列出待补事项。
\begingroup\linespread{1.0}\selectfont\VerbatimInput[fontsize=\fontsize{8.5}{10.5}\selectfont,breaklines=true,breakanywhere=true,numbers=left,numbersep=4pt,xleftmargin=17pt,tabsize=4]{\RepoRoot/prepare_submission.py}\endgroup

\end{document}
